\documentclass[11pt,a4paper]{article}
\ifdefined\pdfmapfile\pdfmapfile{+lm.map}\pdfmapfile{+cm.map}\pdfmapfile{+symbols.map}\fi
\usepackage[margin=25mm,headheight=14pt]{geometry}
\usepackage[T1]{fontenc}
\usepackage{lmodern,amsmath,amssymb,amsthm,microtype}
\usepackage{xcolor,fancyhdr}
\definecolor{navy}{RGB}{25,56,78}
\usepackage[colorlinks=true,linkcolor=navy,citecolor=navy,urlcolor=navy]{hyperref}
\pagestyle{fancy}\fancyhf{}
\fancyhead[L]{\small Real negative zeros of a Fox--Wright family}
\fancyhead[R]{\small Research proof draft}
\fancyfoot[C]{\thepage}
\setlength{\emergencystretch}{2em}
\raggedbottom
\allowdisplaybreaks[1]
\hypersetup{pdftitle={Real Negative Zeros of a Fox-Wright Family},pdfauthor={Keio Research Notes}}
\newtheorem{theorem}{Theorem}
\newtheorem{lemma}[theorem]{Lemma}
\newtheorem{corollary}[theorem]{Corollary}
\newcommand{\LP}{\mathcal{LP}^{+}}
\title{\textbf{Real Negative Zeros of a Fox--Wright Family}\\[5pt]\large A continuous-parameter extension of Craven--Csordas Conjecture~3.17}
\author{Keio Research Notes\thanks{AI-assisted mathematical research. The arguments have undergone independent internal checks; external review and historical priority remain to be established.}}
\date{15 November 2026}
\begin{document}
\maketitle
\begin{abstract}
For every real $m\ge1$, the entire function
$f_m(x)=\sum_{k\ge0}\Gamma(mk+1)x^k/[\Gamma((m+1)k+1)k!]$
has only simple negative real zeros. This proves the positive-integer assertion
in Craven--Csordas (2006), Conjecture 3.17, and extends it continuously in the
parameter. The proof represents an associated cosine transform by an endpoint
atom plus a signed measure of smaller total variation. A positive Laplace kernel
for a normalized gamma quotient gives the required strict dominance uniformly
for all $m\ge1$. No numerical calculation enters the proof.
\end{abstract}

\section{Statement}
For real $m\ge1$, define
\[
f_m(x)=\sum_{k=0}^{\infty}
\frac{\Gamma(mk+1)}{\Gamma((m+1)k+1)}\frac{x^k}{k!},\qquad
c_m=\frac{4m^m}{(m+1)^{m+1}},\qquad
A_m=\sqrt{\frac{2m}{m+1}}.
\]
\begin{theorem}\label{thm:main}
For every real $m\ge1$, $f_m$ belongs to $\LP$ and all its zeros are simple.
For each integer $j\ge0$, there is exactly one zero in
\[
\left(-\frac{(j+1)^2\pi^2}{c_m},
      -\frac{j^2\pi^2}{c_m}\right),
\]
and these are all its zeros. In particular,
$\{\Gamma(mk+1)/\Gamma((m+1)k+1)\}_{k\ge0}$ is a multiplier sequence.
\end{theorem}
Here $\LP$ denotes the class of locally uniform limits of real
polynomials whose zeros are all nonpositive, with nonnegative coefficients.
A nonnegative sequence $\{\gamma_k\}$ is a multiplier sequence if
$\sum a_kx^k\mapsto\sum\gamma_k a_kx^k$ preserves real-rootedness
of every real polynomial.

The positive-integer special case is Conjecture 3.17 of
Craven and Csordas~\cite{CC}, available online in April 2005.
The original paper proves the cases $m=1,2$; the theorem above treats
the original full integer range and also a continuous parameter range. When $m=1$, the identity
$f_1(x)=\cosh\sqrt{x}$ gives the assertion directly. We retain this endpoint
in the uniform argument.

\section{An endpoint-dominance lemma}
\begin{lemma}\label{lem:rouche}
Let $A>0$ and let $\nu$ be a finite real signed Borel measure on $[0,1]$
with total variation $B=\|\nu\|_{\mathrm{TV}}<A$. Then the entire function
\[
F(z)=A\cos z+\int_{[0,1]}\cos(sz)\,\nu(ds)
\]
has only simple real zeros, with exactly one zero in every interval
$(j\pi,(j+1)\pi)$ for $j\in\mathbb Z$.
\end{lemma}
\begin{proof}
The integral defines an entire function because the integrand and all its
complex derivatives are uniformly bounded on compact subsets of $\mathbb C$.
Writing $R=F-A\cos$, we have
\[
|R(x+iy)|\le B\cosh|y|.
\]
At every integer $j$, $|R(j\pi)|\le B<A$, and consequently
$(-1)^jF(j\pi)>0$. There is at least one real zero in each indicated interval.

Choose $Y>0$ with $\tanh Y>B/A$. On the vertical sides $x=\pm N\pi$
of the rectangle $[-N\pi,N\pi]+i[-Y,Y]$,
$|A\cos(x+iy)|=A\cosh y>|R(x+iy)|$. On the horizontal sides,
\[
|A\cos(x\pm iY)|\ge A\sinh Y>B\cosh Y\ge|R(x\pm iY)|.
\]
Rouch\'e's theorem therefore counts exactly $2N$ zeros of $F$ in the
rectangle, with multiplicities, because $\cos z$ has exactly $2N$ there.
The $2N$ intervening real intervals already supply $2N$ distinct zeros.
Thus these zeros are simple and no others occur in the rectangle.
Both $N$ and $Y$ can be taken arbitrarily large, which proves the assertion.
\end{proof}

\begin{lemma}\label{lem:laplace}
Let $\mu$ be a finite nonnegative measure on $(0,\infty)$, of total mass
$L<\log2$, and write
\[
h(x)=\int_0^\infty e^{-xt}\,\mu(dt),\qquad
b(x)=e^{L-h(x)}\quad(x\ge0).
\]
Then the entire function
\[
F(z)=\sum_{k=0}^{\infty}\frac{(-1)^k b(k)z^{2k}}{(2k)!}
\]
satisfies the conclusion of Lemma~\ref{lem:rouche}.
\end{lemma}
\begin{proof}
Put $A=e^L$ and let $\mu^{*n}$ denote additive convolution powers. The real
signed measure
\[
\lambda=A\sum_{n=1}^{\infty}\frac{(-1)^n}{n!}\mu^{*n}
\]
is well-defined in total variation, and
\[
\|\lambda\|_{\mathrm{TV}}
\le A\sum_{n=1}^{\infty}\frac{L^n}{n!}
=A(A-1)<A.
\]
Its Laplace transform is
\[
A+\int_0^\infty e^{-xt}\lambda(dt)
=Ae^{-h(x)}=b(x).
\]
Push $\lambda$ forward by $t\mapsto e^{-t/2}$, obtaining a real signed
measure $\nu$ on $[0,1]$ with no larger total variation. Expanding the
cosine and integrating term by term on compact sets gives
\[
F(z)=A\cos z+\int_{[0,1]}\cos(sz)\nu(ds).
\]
Indeed the coefficient of $(-1)^kz^{2k}/(2k)!$ is
$A+\int e^{-kt}\lambda(dt)=b(k)$. Lemma~\ref{lem:rouche} applies.
\end{proof}

\section{The positive gamma-quotient kernel}
For fixed real $m\ge1$, define on $x\ge0$
\[
b_m(x)=\frac{\Gamma(2x+1)\Gamma(mx+1)}
{\Gamma(x+1)\Gamma((m+1)x+1)c_m^x},\qquad
h_m(x)=\log A_m-\log b_m(x).
\]
Thus $b_m(0)=1$. Stirling's formula gives
\[
\lim_{x\to\infty}b_m(x)=A_m,\qquad
\lim_{x\to\infty}h_m(x)=0,\qquad
\lim_{x\to\infty}h'_m(x)=0.
\]
The last limit also follows directly from the standard digamma asymptotic.
For $t>0$ and $a>0$, set
\[
U_t(a)=\frac1{e^{t/a}-1},\qquad
K_m(t)=U_t(1)+U_t(m+1)-U_t(2)-U_t(m).
\]
\begin{lemma}\label{lem:kernel}
For every real $m\ge1$, $K_m(t)\ge0$ for all $t>0$, and
\[
h_m(x)=\int_0^\infty e^{-xt}\frac{K_m(t)}t\,dt\quad(x\ge0).
\]
In particular the representing positive measure has total mass
\[
\int_0^\infty\frac{K_m(t)}t\,dt
=h_m(0)=\log A_m<\log\sqrt2<\log2.
\]
\end{lemma}
\begin{proof}
For $a>0$,
\[
U_t'(a)=\frac1t
\left(\frac{t/(2a)}{\sinh(t/(2a))}\right)^2.
\]
The function $s\mapsto s/\sinh s$ is strictly decreasing on $(0,\infty)$:
its derivative has the sign of $\sinh s-s\cosh s<0$. Thus $U_t'(a)$ is
strictly increasing in $a$, so $U_t$ is strictly convex. It follows that
\[
K_m(t)=\int_0^1\bigl(U_t'(m+s)-U_t'(1+s)\bigr)\,ds\ge0.
\]
It vanishes identically for $m=1$ and is positive for $m>1$.

The standard trigamma integral gives, for every $a>0$ and $x>0$,
\[
a^2\psi'(ax+1)=\int_0^\infty
\frac{t e^{-xt}}{e^{t/a}-1}\,dt.
\]
Differentiating $h_m$ twice and applying this formula yields
\[
h_m''(x)=\int_0^\infty t e^{-xt}K_m(t)\,dt.
\]
Integrating from infinity twice, using the limits already stated and
Tonelli's theorem for the nonnegative integrand, proves
\[
h_m(x)=\int_x^\infty (u-x)h_m''(u)\,du
=\int_0^\infty e^{-xt}\frac{K_m(t)}t\,dt
\quad(x>0).
\]
These integrations are also justified directly by the endpoint estimates
\[
K_m(t)=\frac{t}{12}
\left(\frac12-\frac1{m(m+1)}\right)+O(t^3)
\quad(t\downarrow0),
\]
and exponential decay as $t\to\infty$ for fixed $m>1$.
At $m=1$ all quantities vanish. Letting $x\downarrow0$ by monotone
convergence and using continuity of the gamma quotient at zero gives the
claimed mass formula. The strict upper bound follows from
$1\le A_m<\sqrt2$ for every finite $m\ge1$.
\end{proof}

\section{Completion of the proof}
\begin{proof}[Proof of Theorem~\ref{thm:main}]
Lemma~\ref{lem:kernel} puts $b_m$ in the setting of
Lemma~\ref{lem:laplace}. Moreover $0\le h_m(x)\le h_m(0)=\log A_m$,
so $1\le b_m(x)\le A_m$. Consequently
\[
f_m(x)=\sum_{k=0}^{\infty}\frac{b_m(k)(c_mx)^k}{(2k)!}
\]
is entire, with order $1/2$. For example the coefficient comparison with
$\cosh\sqrt{c_mx}$ gives both the upper growth bound on $\mathbb C$ and
the matching lower bound on the positive real axis. Its even lift is
\[
F_m(z)=f_m(-z^2/c_m)
=\sum_{k=0}^{\infty}\frac{(-1)^k b_m(k)z^{2k}}{(2k)!}.
\]
By Lemma~\ref{lem:laplace}, all zeros of $F_m$ are simple and real,
with exactly one in each $(j\pi,(j+1)\pi)$.
Every zero $x$ of $f_m$ gives a zero $z$ of $F_m$ with $z^2=-c_mx$.
Therefore $x$ is real and negative. Because $f_m(0)=1$, the substitution
has nonzero derivative at every such lift, so all zeros of $f_m$ are simple.
The positive zeros of $F_m$ give exactly the intervals in the theorem.
In fact the representation also yields an explicit narrower enclosure:
if $\rho_j\in(j\pi,(j+1)\pi)$ is its positive zero, then
\[
 \left|\rho_j-(j+\tfrac12)\pi\right|
 \le \arcsin(A_m-1),\qquad j=0,1,2,\ldots .
\]
Indeed at a real zero the variation bound gives
$|\cos\rho_j|\le A_m-1<1$.

Finally, an entire function of order $1/2$ with only negative real zeros,
positive value at zero, and positive coefficients is in $\LP$: its
Hadamard product is a genus-zero product with those negative zeros and a
positive constant factor. The P\'olya--Schur theorem then identifies the
coefficient sequence of its exponential generating function as a multiplier
sequence, proving the last assertion.
\end{proof}

\section{A coefficient-power consequence}
\begin{corollary}\label{cor:powers}
Fix real $m\ge1$ and $\theta\ge0$. If
$\theta\log A_m<\log2$, then
\[
\Phi_{m,\theta}(x)=\sum_{k=0}^{\infty}
\frac{b_m(k)^\theta x^k}{(2k)!}
\]
belongs to $\LP$, and all its zeros are simple. In particular this holds
uniformly for every $m\ge1$ and every $0\le\theta\le2$.
For $m>1$ the permissible sufficient range is
$0\le\theta<\log2/\log A_m$; for $m=1$ all $\theta\ge0$ work.
\end{corollary}
\begin{proof}
Apply Lemma~\ref{lem:laplace} to the measure $\theta\mu_m$, where
$\mu_m(dt)=K_m(t)dt/t$. Its total mass is $\theta\log A_m$,
and its normalized coefficients are exactly $b_m(k)^\theta$.
The same even-lift argument proves the assertion. The uniform range follows
from $A_m<\sqrt2$ for every finite $m\ge1$.
\end{proof}
For instance, the endpoint $\theta=2$ implies that
\[
\left\{\frac{(2k)!}{k!}
\frac{\Gamma(mk+1)^2}{\Gamma((m+1)k+1)^2}\right\}_{k\ge0}
\]
is a multiplier sequence for every real $m\ge1$. The coefficient identity follows from the definition of $b_m$;
multiplication by the positive geometric sequence $c_m^{2k}$
does not affect the multiplier property.
This assertion is not obtained merely by squaring the original multiplier
sequence, since the additional factor $(2k)!/k!$ is present.

\section{Consequences and scope}
For every integer $d\ge1$ and every real $m\ge1$, the Jensen polynomial
\[
 J_{d,m}(x)=\sum_{k=0}^d\binom dk
 \frac{\Gamma(mk+1)}{\Gamma((m+1)k+1)}x^k
\]
has only negative real zeros. More generally the same diagonal multiplier
preserves the real-rootedness of every input polynomial, in every degree.
These conclusions follow from the multiplier assertion of
Theorem~\ref{thm:main}, not from a finite collection of numerical tests.
The reusable sufficient criterion is Lemma~\ref{lem:laplace}, and
Corollary~\ref{cor:powers} supplies a further continuous family.

The proof is analytic and uniform in all real $m\ge1$. It does not infer
real-rootedness from a positive mixture of real-rooted functions; the measure
in Lemma~\ref{lem:laplace} is allowed to be signed, and the essential fact
is its strict total-variation bound relative to the endpoint atom.
The criterion in Lemma~\ref{lem:laplace} is a reusable sufficient condition
for constructing multiplier sequences from small-mass positive Laplace
kernels. Its sufficient bound $L<\log2$ is not asserted to be optimal.
The proof uses classical gamma-function identities, finite signed measures,
Rouch\'e's theorem, Hadamard factorization, and the P\'olya--Schur theorem;
it requires no unverified theorem from a research corpus.

A targeted literature search found no earlier proof of the full assertion.
This is not a proof of historical priority. Sokal's 2022 paper~\cite{Sokal}
provides relevant later context for hypergeometric membership in $\LP$;
its general parameter obstructions must not be replaced by a blanket
claim that positive mixtures preserve real-rootedness.
The present document is a complete analytic proof draft for external review.

\begin{thebibliography}{9}
\bibitem{CC} T. Craven and G. Csordas,
\emph{The Fox--Wright functions and Laguerre multiplier sequences},
J. Math. Anal. Appl. \textbf{314} (2006), 109--125.
Available online 20 April 2005.
\url{https://doi.org/10.1016/j.jmaa.2005.03.058}.
Conjecture 3.17 and equation (3.8).
\bibitem{DLMF} NIST Digital Library of Mathematical Functions,
\S5.9 (integral representations) and \S5.11 (asymptotic expansions),
\url{https://dlmf.nist.gov/5.9},
\url{https://dlmf.nist.gov/5.11}.
\bibitem{Sokal} A. D. Sokal,
\emph{When does a hypergeometric function ${}_pF_q$ belong to the
Laguerre--P\'olya class $LP^+$?},
J. Math. Anal. Appl. \textbf{515} (2022), 126432.
\url{https://arxiv.org/abs/2204.01045}.
\end{thebibliography}
\end{document}
